% p.21 — الباب الحادي عشر في قال يقول (تابع): المنفيات والمؤكدات
\mspage{21}
% Layout: six tall columns (right-to-left), each a stack of boxed blocks of unequal height.
% Every sub-heading box holds the heading AND, beneath it, the underlying (pre-iʿlāl) form,
% e.g. «مجهول بلم / لَمْ يُقْوَلْ»; the regular forms follow outside the box. Row alignment is
% approximate: all six columns happen to run to 21 lines, so rows are aligned 1:1 here; the
% \cline marks reproduce the printed box boundaries inside each column.
% Column 1 continues the مضارع بلم (معروف) block begun on p.20.
\begin{sarf}{|c|c|c|c|c|c|}
لَمْ تَقُوْلَا & لَمْ أُقَلْ & لَنْ يُقَالَا & لَتَقُوْلَنَّ & لَتُقَلْنَانِّ & \thead{أمر حاضر بالنون}\\
لَمْ تَقُوْلُوْا & لَمْ نُقَلْ & لَنْ يُقَالُوْا & لَتَقُوْلَانِّ & لَأُقَالَنَّ & اُقْوُلَنَّ\\\cline{2-2}\cline{6-6}
لَمْ تَقُوْلِي & \thead{مضارع بلن} & لَنْ تُقَالَ & لَتَقُوْلُنَّ & لَنُقَالَنَّ & قُوْلَنَّ\\\cline{5-5}
لَمْ تَقُوْلَا & لَنْ يَقْوُلَ & لَنْ تُقَالَا & لَتَقُوْلِنَّ & \thead{منفي بالنون} & قُوْلَانِّ\\\cline{2-2}
لَمْ تَقُلْنَ & لَنْ يَقُوْلَ & لَنْ يُقَلْنَ & لَتَقُوْلَانِّ & لَا يَقْوُلَنَّ & قُوْلُنَّ\\\cline{5-5}
لَمْ أَقُلْ & لَنْ يَقُوْلَا & لَنْ تُقَالَ & لَتَقُلْنَانِّ & لَا يَقُوْلَنَّ & قُوْلِنَّ\\\cline{5-5}
لَمْ نَقُلْ & لَنْ يَقُوْلُوْا & لَنْ تُقَالَا & \com{لَا قُوْلَنَّ}{كذا في الأصل، بألف بلا همزة؛ ولعل الصواب: لَأَقُوْلَنَّ} & \thead{مجهول بالنون} & قُوْلَانِّ\\\cline{1-1}
\thead{مجهول بلم} & لَنْ تَقُوْلَ & لَنْ تُقَالُوْا & لَنَقُوْلَنَّ & لَا يُقْوَلَنَّ & قُلْنَانِّ\\\cline{4-4}\cline{5-5}\cline{6-6}
لَمْ يُقْوَلْ & لَنْ تَقُوْلَا & لَنْ تُقَالِي & \thead{مجهول بالنون} & لَا يُقَالَنَّ & \thead{نهي حاضر بالنون}\\\cline{1-1}\cline{5-5}
لَمْ يُقَلْ & لَنْ يَقُلْنَ & لَنْ تُقَالَا & لَيُقْوَلَنَّ & \thead{أمر غائب بالنون} & لَا تَقْوُلَنَّ\\\cline{4-4}\cline{6-6}
لَمْ يُقَالَا & لَنْ تَقُوْلَ & لَنْ تُقَلْنَ & لَيُقَالَنَّ & لِيَقْوُلَنَّ & لَا تَقُوْلَنَّ\\\cline{5-5}
لَمْ يُقَالُوْا & لَنْ تَقُوْلَا & لَنْ أُقَالَ & لَيُقَالَانِّ & لِيَقُوْلَنَّ & لَا تَقُوْلَانِّ\\\cline{5-5}
لَمْ تُقَلْ & لَنْ تَقُوْلُوْا & لَنْ نُقَالَ & لَيُقَالُنَّ & \thead{مجهول بالنون} & لَا تَقُوْلُنَّ\\\cline{3-3}
لَمْ تُقَالَا & لَنْ تَقُوْلِي & \thead{مضارع بالنون} & لَتُقَالَنَّ & لِيُقْوَلَنَّ & لَا تَقُوْلِنَّ\\\cline{5-5}
لَمْ يُقَلْنَ & لَنْ تَقُوْلَا & لَيَقْوُلَنَّ & لَتُقَالَانِّ & لِيُقَالَنَّ & لَا تَقُوْلَانِّ\\\cline{3-3}\cline{5-5}
لَمْ تُقَلْ & لَنْ تَقُلْنَ & لَيَقُوْلَنَّ & لَيُقَلْنَانِّ & \thead{نهي غائب بالنون} & لَا تَقُلْنَانِّ\\\cline{6-6}
لَمْ تُقَالَا & لَنْ أَقُوْلَ & لَيَقُوْلَانِّ & لَتُقَالَنَّ & لَا يَقْوُلَنَّ & \thead{تمرين}\\\cline{5-5}\cline{6-6}
لَمْ تُقَالُوْا & لَنْ نَقُوْلَ & لَيَقُوْلُنَّ & لَتُقَالَانِّ & لَا يَقُوْلَنَّ & جَالَ\\\cline{2-2}\cline{5-5}
لَمْ تُقَالِي & \thead{مجهول بلن} & لَتَقُوْلَنَّ & لَتُقَالُنَّ & \thead{مجهول بالنون} & جَازَ\\
لَمْ تُقَالَا & لَنْ يُقْوَلَ & لَتَقُوْلَانِّ & لَتُقَالِنَّ & لَا يُقْوَلَنَّ & جَاعَ\\\cline{2-2}\cline{5-5}
لَمْ تُقَلْنَ & لَنْ يُقَالَ & لَيَقُلْنَانِّ & لَتُقَالَانِّ & لَا يُقَالَنَّ & حَاجَ\\
\end{sarf}
